%
En gass er i tilstand A med trykk $p_\text{A} = \SI{100}{\kilo\pascal}$, volum $V_\text{A} = \SI{4.0}{\litre}$ og temperatur $T_\text{A} = \SI{300}{\kelvin}$. Gassen blir først presset sammen isotermt til volumet \SI{1.0}{\litre} (tilstand B). Deretter blir den varmet opp ved konstant volum til \SI{600}{\kelvin} (tilstand C).
\begin{enumerate}[a)]
    \item Finn trykket i tilstand B.
    \item Finn trykket i tilstand C.
    \item Skisser de to prosessene i et $pV$-diagram.
    \item Gjør gassen noe arbeid fra B til C? Begrunn svaret.
\end{enumerate}
%
